% SPDX-License-Identifier: LGPL-3.0-only
\chapter{Implementation}
\label{ch:implementation}

\section{A Reference Implementation}
Listing~\ref{lst:midpoint} implements Equation~\eqref{eq:midpoint}.
It rejects nonpositive subinterval counts before computing the step size.

\begin{lstlisting}[language=Python,caption={Composite midpoint rule.},label={lst:midpoint}]
def midpoint(f, a, b, n):
    if n < 1:
        raise ValueError("n must be positive")
    h = (b - a) / n
    return h * sum(f(a + (k + 0.5) * h) for k in range(n))
\end{lstlisting}

\section{Checking the Result}
Theorem~\ref{thm:affine} provides an exact test case.
Table~\ref{tab:convergence} instead uses \(f(x)=x^2\) to show how the
approximation changes when the number of subintervals doubles.

\begin{table}[htbp]
  \centering
  \caption{Midpoint estimates for the integral of \(x^2\) on \([0,1]\).}
  \label{tab:convergence}
  \begin{tabular}{rrr}
    \toprule
    \(n\) & Estimate & Absolute error \\
    \midrule
    1 & 0.250000 & 0.083333 \\
    2 & 0.312500 & 0.020833 \\
    4 & 0.328125 & 0.005208 \\
    \bottomrule
  \end{tabular}
\end{table}

\begin{remark}
These values check the implementation for one integrand. They do not
establish an error bound for arbitrary continuous functions.
\end{remark}
